\documentclass[11pt]{article}
\usepackage{fontspec}
\usepackage[french,english]{babel}
\usepackage{amsmath}
\usepackage{amssymb}
\usepackage{amsthm}
\usepackage{mathrsfs}
\usepackage{tikz-cd}
\usepackage{enumitem}
\usetikzlibrary{babel}

\newcommand{\printedpage}[1]{\marginpar{\footnotesize p.~#1}}
\newcommand{\msp}[1]{\marginpar{\footnotesize [#1]}}
\newcommand{\Ob}{\operatorname{Ob}}
\newcommand{\Fl}{\operatorname{Fl}}
\newcommand{\card}{\operatorname{card}}
\newcommand{\Ind}{\operatorname{Ind}}
\newcommand{\Hom}{\operatorname{Hom}}
\newcommand{\id}{\mathrm{id}}
\newcommand{\limind}{\varinjlim}
\newcommand{\cercle}[1]{\tikz[baseline=(c.base)]\node[draw,circle,inner sep=1pt](c){#1};}
\newcommand{\defeq}{\overset{\text{def}}{=}}

\begin{document}

\printedpage{134}
% head: A. Grothendieck : Dérivateurs, Ch. XVIII, Édition des Universités de Montpellier II et Paris VII
% foot: Version 17 décembre 2025
\begin{itemize}[leftmargin=2.5em]
\item[2)] The direct limits large relative to $\pi_0$ is of cardinality $\leqslant \pi$ exist in $I$ (follows from $\mathrm{L}_{\pi_0}$ on $\mathcal{M}$ (cf.\,\cercle{d})). % [sic] est de cardinal -> et de cardinal, coordination ("grandes devant π0 et de cardinal ≤ π")
\item[3)] The functor $\xi : I \longrightarrow \mathcal{M}$ commutes with them.
\item[4)] The $x_i$ are in $\mathcal{M}_\pi$.
\end{itemize}
Since $f$ commutes with the $\limind$ large relative to $\pi_0$ (cf.\,\cercle{1}), a fortiori with those large relative to $\pi$, we therefore have
\[ f(x) \;\simeq\; \varinjlim_{I} f(x_i) \ , \]
I say that the direct system $(f(x_i))_{i\in I}$ in $\mathcal{N}$ is $\pi_0$-$\pi$-adapted. Conditions 1), 2) are \textit{ne varietur} (concern only $I$), the functor $I \longrightarrow \mathcal{N}$ is the composite $I \xrightarrow{\xi} \mathcal{M} \xrightarrow{f} \mathcal{N}$, hence commutes with the $\limind$ large relative to $\pi_0$ and of cardinality $\leqslant \pi$ (\footnote[155]{this commutation should have a name, condition $\mathrm{L}(\pi_0,\pi)$ on a functor.}), since this is so for $\xi$ and $f$, finally the $f(x_i)$ are in $\mathcal{N}_\pi$ since by \cercle{c} $f(\mathcal{M}_\pi) \subset \mathcal{N}_\pi$.

\msp{page 353}

Moreover $f(x) \in \Ob \underline{T}$ (since $x \in \Ob \underline{S}$, where $\underline{S} \defeq f^{-1}(\underline{T})$). Thus we can apply 4.13.22 a $\Longrightarrow$ c to the situation $(\mathcal{N}, \pi_0, \pi, \underline{S})$ and to the direct system $(f(x_i))_{i\in I}$, to conclude that the subcategory $I'$ of $I$ consisting of the $i$ such that $f(x_i) \in \Ob \underline{S}$ \texttt{[rather Ob}~$\underline{T}$\texttt{]}, i.e. such that $x_i \in \underline{S}$, hence $x_i \in \underline{S}(\pi)$, is \textit{cofinal} in $I$. Thus we have
\[ x \;\simeq\; \varinjlim_{i\in I'} x_i \ , \]
with this time the $x_i$ in $\underline{S}(\pi)$, and $I'$ large relative to $\pi$ (since cofinal and full in $I$, which is large relative to $\pi$), which proves that $x$ is in the essential image of $\Ind_\pi(\underline{S}(\pi))$, q.e.d.

\textbf{Question 4.13.33.} We have just seen that if $f : \mathcal{M} \longrightarrow \mathcal{N}$ is an accessible functor between accessible categories, then the inverse image of an accessible subset of $\Ob \mathcal{N}$ is an accessible subset of $\Ob \mathcal{M}$. I conjecture that the similar statement, not for direct images, but for \textit{essential images}, is
\msp{page 354}
true, i.e. that for every subcategory $\underline{S}$ of $\mathcal{M}$ accessible in $\mathcal{M}$, its essential image in $\mathcal{N}$ is accessible in $\mathcal{N}$ (\footnote[156]{False, even if we take the $\natural$-closure of the essential image (``$\natural$-essential image'').}). In this question, we can assume that $\underline{S} = \mathcal{M}$, and the question comes down to this: does every accessible functor $f : \mathcal{M} \longrightarrow \mathcal{N}$ between accessible categories factor into an (accessible) \textit{essentially surjective} functor, and an accessible fully faithful functor $\mathcal{M} \longrightarrow \underline{T} \longrightarrow \mathcal{N}$ (where $\underline{T}$ is itself accessible) (\footnote[157]{We must at least restrict ourselves to essentially $\natural$-surjective – but even that is too optimistic.})~? If so, this would seem to me perhaps the deepest result (perhaps the only truly deep one) in category theory. % [sic] vraie -> vrai, agreement with « l'énoncé similaire »

Here is a question that seems to me very close. Let $\mathcal{M}$ be an accessible category, $A$ a small full subcategories. Suppose that $\mathcal{M}$ satisfies $\mathrm{L}_{\pi_0}$ (e.g. $\pi_0$ useful for $\mathcal{M}$), and let $\underline{S}$ be the smallest strictly full subcategory of $\mathcal{M}$ containing $A$, and stable under limits large relative to $\pi_0$. I conjecture that $\underline{S}$ is accessible in $\mathcal{M}$. % [sic] sous-catégories -> sous-catégorie, singular after « une »

\msp{page 355}

Likewise, if $\Phi$ is a small set of small categories, such that $\mathcal{M}$ is stable under $\limind$ of type $\Phi$, we can take the envelope $\underline{S}$ of $A$ for the $\limind$ large relative to $\pi_0$ \textit{and} the
\printedpage{135}
% head: A. Grothendieck : Dérivateurs, Ch. XVIII, Édition des Universités de Montpellier II et Paris VII
% foot: Version 17 décembre 2025
$\limind$ of type $\Phi$, and I again conjecture that $\underline{S}$ is accessible. In this direction, we have at least the following very partial result.

\textbf{Proposition 4.13.34.} \textit{Let $\mathcal{M}$ be an accessible category stable under small filtering $\limind$, and let $\underline{S}$ be a strictly full subcategory of $\mathcal{M}$. Equivalent conditions:}
\begin{itemize}[leftmargin=2.5em]
\item[\textit{a)}] \textit{$\underline{S}$ is accessible in $\mathcal{M}$, and stable in $\mathcal{M}$ under small filtering $\limind$.}
\item[\textit{b)}] \textit{There exists a small full subcategory $A$ of $\mathcal{M}$, such that $\underline{S}$ is the smallest strictly full subcategory of $\mathcal{M}$ containing $A$ and stable under small filtering $\limind$.}
\end{itemize}

\msp{page 356}

We have a) $\Longrightarrow$ b), for if $\underline{S}$ is $\pi$-accessible in $\mathcal{M}$, it is immediate (without any other stability condition on $\underline{S}$) that $\underline{S}$ is the smallest strictly full subcategory of $\mathcal{M}$ containing $\underline{S}(\pi) \defeq \underline{S}\cap\mathcal{M}_\pi$ and stable under $\limind$ large relative to $\pi$. If then we already know that $\underline{S}$ is even stable under \textit{all} small filtering $\limind$, it follows that $\underline{S}$ is the smallest subcategory of $\mathcal{M}$ having this stronger stability property, and containing $A$.

Let us prove b $\Longrightarrow$ a. Let $\pi$ be a useful cardinal for $\mathcal{M}$, and such that $A \subset \mathcal{M}_\pi$. Let $B$ be the smallest strictly full subcategory of $\mathcal{M}$ containing $A$ and stable under filtering $\limind$ of cardinality $\leqslant \pi$. Thus we have $B \hookrightarrow \mathcal{M}_\pi$ and $B \subset \underline{S}$, and $\underline{S}$ is also the L-envelope of $B$. Besides $B$ is stable under
\msp{page 357}
direct factors. Let us then consider the essential image $\underline{S}'$ of $\Ind_\pi(B)$ in $\mathcal{M}$. We have $\underline{S}' \subset \underline{S}$, and everything comes down to proving that $\underline{S}' = \underline{S}$, from which it will indeed follow that $\underline{S}$ is $\pi$-accessible in $\mathcal{M}$, a fortiori accessible in $\mathcal{M}$, whence a). To prove that $\underline{S}' = \underline{S}$, by definition of $\underline{S}$ as the L-envelope of $A$, it suffices to see that $\underline{S}'$ is stable under small filtering $\limind$. Now this follows from the fact that $B \subset \mathcal{M}_\pi$ is stable in $\mathcal{M}_\pi$ under filtering $\limind$ of cardinality $\leqslant \pi$, and from 4.13.19, which implies at least, since $B = \underline{S}' \cap \mathcal{M}_\pi$ (by virtue of 4.13.4) that $\underline{S}'$ is stable under filtering $\limind$ \textit{of cardinality} $\leqslant \pi$. That it is even stable under arbitrary small filtering $\limind$ follows from this, and from the fact that it is so under $\limind$ large relative to $\pi$. (Cf. 4.13.21 \dots) % [sic] filtrante quelconques -> filtrantes quelconques, plural agreement

\msp{page 358}

\textbf{Definition 4.13.35.} Let $\mathcal{M}$ be an accessible category, stable under small filtering $\limind$, $S$ a subset of $\Ob \mathcal{M}$, $\underline{S}$ the corresponding full subcategory in $\mathcal{M}$. We say that $S$ is \textit{L-accessible}, or that $\underline{S}$ is L-accessible in $\mathcal{M}$, if $\underline{S}$ is accessible in $\mathcal{M}$, and in addition stable under small filtering $\limind$ (\footnote[158]{\textbf{N.B.} The set of L-accessible subsets of $\Ob \mathcal{M}$ is stable under small intersections, and under \textit{finite} unions. The inverse image of an L-accessible subset under a functor (\footnotemark[159]) $f : \mathcal{M} \longrightarrow \mathcal{N}$ of accessible categories is L-accessible.}).\footnotetext[159]{an \textit{$L_0$-functor}: which commutes with small filtering $\limind$.}

\textbf{Proposition 4.13.36.} \textit{Let $\mathcal{M}$, $S$, $\underline{S}$ be as in 4.13.35, and let $\pi$ be a useful cardinal for $\mathcal{M}$. Suppose $\underline{S}$ stable in $\mathcal{M}$ under small filtering $\limind$, or only $\underline{S}(\pi) \defeq \underline{S} \cap \mathcal{M}_\pi$ stable under small filtering $\limind$ of cardinality $\leqslant \pi$. Equivalent conditions:}
\begin{itemize}[leftmargin=2.5em]
\item[\textit{a)}] \textit{$\underline{S}$ is $\pi$-accessible. (It will automatically be stable under small filtering $\limind$, hence L-accessible.)}
\item[\textit{b)}] \textit{For every $x \in \Ob \mathcal{M}$, we have $x \in \Ob \underline{S}$ if and only if $\underline{S}(\pi)/x$ is cofinal in $\mathcal{M}_\pi/x$.}
\printedpage{136}
% head: A. Grothendieck : Dérivateurs, Ch. XVIII, Édition des Universités de Montpellier II et Paris VII
% foot: Version 17 décembre 2025
\item[\textit{c)}] \textit{For every $x \in \Ob \mathcal{M}$, and every representation $x \simeq \varinjlim\limits_{i\in I} x_i$}
\msp{page 359}
\textit{which is L-$\pi$-adapted (cf. page 343, c), we have this: $x \in \Ob S$ if and only if $I_{\underline{S}}$ is cofinal in $I$.}
\item[\textit{c}$'$\textit{)}] \textit{As in c), with $I$ ordered.}
\item[\textit{c}$''$\textit{)}] \textit{(If we have chosen, for every $x \in \Ob \mathcal{M}$,} one \textit{$\pi$-adapted representation $x \simeq \varinjlim\limits_{I} x_i$): $x \in \Ob \underline{S} \Longleftrightarrow$ the set $I_{\underline{S}}$ is cofinal in $I$.}
\end{itemize}
\textsc{Proof}. This is essentially the statement of 4.13.22 in the special case where $\Phi =$ set of \textit{filtering} categories of cardinality $\leqslant \pi$, but where, to simplify, we have assumed beforehand that $\underline{S}(\pi)$ satisfies $\mathrm{L}(\Phi)$. The equivalence of a) b) c) is contained in 4.13.22, as is c $\Longrightarrow$ c$'$. But for c$'$ $\Longrightarrow$ a, the auxiliary hypothesis $\pi^{\pi_0} = \pi$ and that the $J \in \Phi$ will be large relative to $\pi_0$, is not satisfied here, and we need an independent proof, modeled on that of pages 333-336. For this it suffices, at the place \texttt{[of]} $\square$, page 333, to invoke the % [sic] un démonstration -> une démonstration, gender agreement

\textbf{Lemma 4.13.37.} \textit{Let $\mathcal{M}$ be a category}
\msp{page 360}
\textit{that is accessible and stable under filtering $\limind$, $\pi$ a useful cardinal for $\mathcal{M}$. Then for every $x \in \Ob \mathcal{M}$, we have a representation}
\[ x \;=\; \varinjlim_{i\in I} x_i \]
\textit{L-$\pi$-adapted (i.e. with}
\begin{itemize}[leftmargin=2.5em]
\item[\textit{1)}] \textit{$I$ equivalent to a small category, and $I$ large relative to $\pi$,}
\item[\textit{2)}] \textit{$I$ stable under filtering $\limind$ of cardinality $\leqslant \pi$,}
\item[\textit{3)}] \textit{$I \longrightarrow \mathcal{M}$ commutes with these $\limind$ ,}
\item[\textit{d)}] \textit{$x_i \in \Ob \mathcal{M}_\pi \ \forall\, i \in I$),} % [sic] d) -> 4), the items are numbered 1), 2), 3)
\end{itemize}
\textit{with in addition $I$} ordered.

Indeed, we start from
\[ x \;=\; \varinjlim_{j\in J} x_j \ , \qquad\quad J \;\simeq\; \underbrace{(\mathcal{M}_\pi/x)_0}_{\substack{\text{reduced subcategory}\\ \text{equivalent to } \mathcal{M}_\pi/x}} \qquad . \]
We define $I$ as the set of \textit{filtering} subcategories $J'$ of $J$, of cardinality \texttt{[``cardinality'' crossed out]} $\leqslant \pi$, $I$ ordered by inclusion, and we set, for $J' \in I$,
\[ x_{J'} \;=\; \varinjlim_{j\in J'} x_j \ . \]
It is immediate that we thus obtain a direct system in $\mathcal{M}$, indexed by $I$.

\msp{page 361}

On the other hand, it is clear that $I$ is small, I say that $I$ is stable under filtering $\limind$ of cardinality $\leqslant \pi$, and that the functor $J' \longmapsto x_{J'} : I \longrightarrow \mathcal{M}$ commutes with them, finally that the $x_{J'}$ are in $\mathcal{M}_\pi$. It remains to see that $I$ is large relative to $\pi$, i.e. that if we have $(J'_r)_{r\in A}$, with $\card A \leqslant \pi$, the $J'_r$ in $I$, then $\exists\, J' \in I$ which dominates the $J'_r$. Let $K$ be the subcategory of $J$ generated by the $J'_r$, we easily see that $\card K \leqslant \pi$. Everything \texttt{[``everything'' crossed out]} comes down to proving:

\printedpage{137}
% head: A. Grothendieck : Dérivateurs, Ch. XVIII, Édition des Universités de Montpellier II et Paris VII
% foot: Version 17 décembre 2025
\textbf{Lemma 4.13.38} (\footnote[160]{\textbf{N.B.} For the proof of 4.13.37, it suffices to prove 4.13.38 in the case where $J$ is ordered, in which case the contortions at the end are unnecessary.}). \textit{Let $J$ be a category large relative to $\pi$, $K$ a subcategory of $J$ of cardinality $\leqslant \pi$. Then there exists a subcategory $K'$ that is} filtering \textit{in $J$, containing $K$, of cardinality $\leqslant \pi$. If we assume that we are not, in the case where $J$ contains a maximal object, and that $K$ contains all the maximal objects of $J$,}
\msp{page 362}
\textit{we can even take $K'$ such that $\Ob K' = \Ob K \cup \{e\}$, where $e$ is a} final object \textit{of $K'$, and $K$ is a full subcategory of $K'$. (Thus if $K' \neq K$, $K'$ is isomorphic to the right cone constructed on $K$, ``obtained from $K$ by adding a final object to it''.)}

\textsc{Proof}. If $K$ has a final object, we take $K' = K$. Otherwise, since $J$ is large relative to $\card \Ob K$, there exists a $j \in \Ob J$ which dominates all the objects of $K$. Let us choose, for every $k \in \Ob K$, an arrow
\[ \alpha_k \;:\; k \longrightarrow j \ . \]
For every arrow
\[ u \;:\; k' \longrightarrow k \]
in $K$, let us consider the double arrow
\[ (\alpha_k \circ u, \alpha_{k'}) \;:\; k' \rightrightarrows j \ . \]
The set of these double arrows, with target $j$, is of cardinality $\leqslant \card \Fl K \leqslant \pi$, so $J$ being large relative to $\pi$, there exists an arrow $\rho : j \longrightarrow e$ in $J$ which equalizes all these double arrows. Let $K' \subset J$, having as objects $\Ob K \cup \{e\}$, and having as arrows, in addition to
\msp{page 363}
those of $K$, the arrows $\rho\alpha_k : k \longrightarrow e$ (for $k \in \Ob K$), and $\id_e$. \textit{If} $e \notin \Ob K$, we see that this is indeed a subcategory, admitting $e$ as final object and $K$ as full subcategory. If we had $e \in \Ob K$, and if there exists an object $e'$ of $J$ which dominates $e$ and which is not in $K$, we immediately reduce to the preceding case, replacing $e$ by $e'$.

$\Biggl($ (\footnote[161]{The construction which follows is contained in the more general and more precise statement 4.13.39.}). Thus there remains the case where there exists an object $e$ in $K$ which is a greatest element in $K$, and such that every object of $J$ dominating $e$ is in $K$. This implies, $J$ being filtering, that $e$ is a greatest object of $J$. (For if $x \in J$, $\exists\, x'$ a common upper bound for $x$ and $e$, hence $x'$ is in $K$, hence dominated by $e$, which therefore dominates $x$.) We see as before that there exists an arrow $p : e \longrightarrow j$ in $J$, which equalizes all the double arrows $k \rightrightarrows e$ with target $e$ \textit{in}~$K$. Composing with $j \longrightarrow e$ if necessary, we can assume $j = e$. Thus we have an endomorphism $p$ of $e$ in $J$, such that $pu = pv$ for every double arrow $k \overset{u}{\underset{v}{\rightrightarrows}} e$ in $K$, with target $e$.

\msp{page 364}

For $k \in \Ob K$, let
\[ f_k \;:\; k \longrightarrow e \]
be the common value of the $pu$, for $u \in \Hom_K(k,e)$. We see at once that if we have $k' \xrightarrow{v} k$ in $K$, then
\[ f_{k'} \;=\; f_k \circ v \ . \]
\printedpage{138}
% head: A. Grothendieck : Dérivateurs, Ch. XVIII, Édition des Universités de Montpellier II et Paris VII
% foot: Version 17 décembre 2025
In particular, taking $k' = e$, we find
\[ f_k \circ v \;=\; f_e \,(=\ p\ ) \qquad\quad \text{if } v : e \longrightarrow k \ . \]
Let then $K'$ be the subcategory of $J$ having the same objects as $K$, and as arrows, in addition to those of $K$, the composite arrows $k \xrightarrow{u} e \xrightarrow{p^r} e \xrightarrow{v} k'$, with $u$, $v$ in $K$, and $r \geqslant 1$, i.e. the composite arrows $k \xrightarrow{f_k} e \xrightarrow{p^{r'}} e \xrightarrow{v} k'$, with $k, k' \in \Ob K$, $v \in \Fl K$, $r' \geqslant 0$. We must see that it is stable under composition (to indeed have a subcategory). This is clear for $vu$ if $u$ or $v$ are in $\Fl K$. If \texttt{[neither]} $u$ nor $v$ is, we must take a composite
\[
\begin{tikzcd}[column sep=large]
k \arrow[r, "u_1"] & e \arrow[r, "p^r", "r\geqslant 1"'] & e \arrow[r, "u_2"] \arrow[rr, bend right=20] & k' \arrow[r, "v_1"] & e \arrow[r, "p^s", "s\geqslant 1"'] & e \arrow[r, "v_2"] & k'' \ ,
\end{tikzcd}
\]
with $u_1$, $u_2$, $v_1$, $v_2$ in $\Fl K$, $r, s \geqslant 1$. Now we have
\[ p(v_1u_2) \;=\; p \ , \qquad\quad \text{whence} \qquad\quad p^s(v_1u_2) \;=\; p^{s-1}(pv_1u_2) \;=\; p^{s-1}p \;=\; p^s \ , \]
so the composite is
\[ k \xrightarrow{\;u_1\;} e \xrightarrow{\;p^{r+s}\;} e \xrightarrow{\;v_2\;} k'' \ , \]
which
\msp{page 365}
is therefore indeed in $\Fl K'$.

Thus, $K'$ is a subcategory of $J$, but is it filtering? Condition PS 1$'$ is obvious, since $e$ is a greatest element of $J$, there remains condition PS 2 of the double arrow. We easily see that this condition is satisfied in $K'$, if and only if it is satisfied for every double arrow
\[ e \overset{p^r}{\underset{p^s}{\rightrightarrows}} e \qquad\quad \text{with } r, s \geqslant 0 \ . \]
(Indeed, we immediately reduce to the double arrows with target $e$, and note that every arrow of $K'$ with target $e$ can be written as a composite $k \xrightarrow{u} e \xrightarrow{p^r} e \xrightarrow{v} e$, with $r \geqslant 0$ and $u$, $v$ arrows in $K'$. Composing with $e \xrightarrow{p} e$, we then find, since $pv = pu$, $k \xrightarrow{u} e \xrightarrow{p^{r+1}} e$, that is $k \xrightarrow{f_k} e \xrightarrow{p^r} e$.) But we see at once that the possibility of equalizing $(p^r, p^s)$ in $K'$ is equivalent to the existence of $t \geqslant 0$ such that $p^{r+t} = p^{s+t}$. Thus $K'$ is filtering if and only if $r \longmapsto p^r$ is a stationary sequence (take $r = 0$, $s = 1$, whence $p^r = p^{r+1}$ for $r$ large). This is
\msp{page 366}
not necessarily the case.

Let us therefore set $K = K_0$, $K' = K_1$, and construct by induction an increasing sequence of subcategories $K_i$ ($i \in \mathbf{N}$), with $\Ob K_i = \Ob J$, and (for $i \geqslant 1$) $K_i$ being obtained from $K_{i-1}$ by adding to it an arrow $p_i : e \longrightarrow e$, as above. Thus every double arrow in a $K_i$ can be equalized in $K_{i+1}$. Let us take $K_\infty = \limind K_i$, $K_\infty$ is a \textit{filtering} subcategory of $J$, containing $K$ and of cardinality $\leqslant \pi$, q.e.d. $\Biggr)$

We can make this construction more precise as follows:

\end{document}
